\section{Limitations and Future Work}
\label{sec:limitations}

This study presents the first quantitative evaluation of re-identifying a crack in a single visit using a geometric reference.
The results are limited by the small number of clustered walls, single-annotator labels, and partial registration coverage.
Each item pairs a finding with a verification step.
\begin{enumerate}
\item \textit{Scale to more walls and visits.}
Collect data from more walls and surface types and annotations from additional annotators, tightening wall-clustered confidence intervals and making rankings more reliable; report leave-one-wall-out results over all 17 walls instead of the 6/11 split.
Most importantly, add a second session with different distances and lighting for a genuine cross-visit benchmark.
\item \textit{Add provenance-free annotation.}
Labels come from a single annotator (Section~\ref{sec:annotation}) and were built by homography propagation, so the geometric advantage may depend on provenance. To test whether it persists independently of label provenance, create labels that do not rely on homography propagation and release per-pair score matrices for wall-clustered confidence intervals and threshold-free metrics (AUROC, AP).
\item \textit{Close the coverage gap.}
The geometric reference covers 48\% of valid query--gallery pairs (15,328 of 31,822), from the 97 of 301 same-wall image pairs (32\%) that register. It tends to fail on smooth plaster, with extent unquantified.
Surface-adaptive registration (learned local features, direct photometric alignment, or sparse depth from consumer devices) would allow a more thorough evaluation and better separate geometric evidence from homography-based provenance.
\item \textit{Learn from geometric supervision.}
The geometric reference is untrained and not deployed.
Location cues could be obtained without per-pair registration at inference time via registration-derived pseudo-labels or fine-tuning a foundation model on homography correspondences.
\item \textit{Assess the ticket lifecycle.}
Define task metrics such as duplicate rate and missed-new-defect rate, and test segmentation, re-identification, and joint ticket assignment in a closed-loop simulation.
\end{enumerate}
Together, these steps would turn the within-visit reference into a multi-visit component usable in operational ticket systems.
